\scanpage{93}
% archive sheet number: 95

\uline{Courbes elliptiques réelles} \deleted{Elles sont} La catégorie \unsure{de|des} ces courbes (avec morphismes pas néc. des iso) revient à celle des \deleted{\unsure{$\pi$|$p$}} systèmes $(\pi,\sigma,c)$, où $\pi$ est un $\mathbb Z$-module libre, \inserted{de rg $2$} \deleted{$\sigma$ une} $c$ une structure \deleted{conf} conforme orientée sur $\pi_{\mathbb R}=\pi\otimes_{\mathbb Z}\mathbb R$, et $\sigma$ un anti-automorphisme du syst. $(\pi,c)$, qui peut s'interpréter comme une involution de $\pi$, renversant l'\uline{orientation} \inserted{de $\pi\otimes_{\mathbb Z}\mathbb R$ i.e. de déterminant $-1$} (donc $\ne\pm\mathrm{id}_\pi$), et conservant la structure conforme.

Pour le couple $(\pi,\sigma)$ il y a exactement deux possibilités, à iso. près. \deleted{Pour} (elles sont d'ailleurs isogènes...). Il reste, pour chacune d'elles, à trouver les structures conformes qui sont invariantes par $\sigma$. On peut les normaliser par \deleted{$\unsure{\sigma Q=Q|\sigma Q=\sigma}$}

\[
Q(u_+)\deleted{(=Q(v))}=1
\]

\margin{\uline{NB}\quad $u_+$ est une base de $\pi_+$\\ $u_-$ \rule{6em}{0.4pt} $\pi_-$}

donc ce sont les structures de la forme

\[
Q(xu_++yu_-)=x^2+ay^2
\]

où $a\in\mathbb R^{*+}$. Il y a donc deux familles de courbes elliptiques réelles, paramétrées chacune par $\mathbb R^{*+}$, on peut les désigner par $E_a,E'_a$ ($a\in\mathbb R^{*+}$). \deleted{Considérons}

Considérons les autom \deleted{de ces cour} \inserted{d'une courbe elliptique} définie par $(\pi,\sigma,c)$, ils correspondent aux \deleted{\unsure{un tel doit|un tel est}} autom $\theta$ de $\pi$ qui a) commutent à $\sigma$ et b) respectent la structure conforme \inserted{orientée} $c$, \deleted{\unsure{en particulier|ce qui peut s'interpréter|ou plutôt}} ou ce qui revient au $=$ (étant d'ordre fini) la forme $Q$ \inserted{(en plus de l'orientation)}. \deleted{Donc} $\theta$ transforme $\pi^{+}_{-}$ en $\pi^{+}$, \deleted{Montrons la symétrie centrale} \deleted{\unsure{fixe|ainsi que}} \inserted{il transforme la $\unsure{2^{\text{de}}|1^{\text{ère}}}$} en un multiple \inserted{soit} $\pm e_0$, $\unsure{\mathrm{OPS}|\mathrm{DPS}}$ quitte à composer avec la symétrie \unsure{centrale $-1$)|(\mbox{i.e.} de $-1$)} que c'est $e_0$. Mais les seul autom de $\mathbb R e_0+\mathbb R e_1$ qui respecte $Q$, qui fixe $e_0$ et qui n'est pas l'identité, est $\sigma$, \deleted{\unsure{qui inverse|précisément}} \inserted{or il} inverse l'orientation ! Donc il n'y a que l'autom universel, correspondant à $-\mathrm{id}_{\unsure{\pi|V}}$.
